% Chapter 4: Industry Drivers & Trends + PESTLE
\chapter{Drivers, Tendencias y Análisis PESTLE}

\subsection{Drivers Primarios de Crecimiento}

\subsubsection{Transición Energética y Electrificación}
La descarbonización global es el principal motor de demanda de minerales críticos. La adopción masiva de vehículos eléctricos, la expansión de redes eléctricas y la construcción de infraestructura de energías renovables (eólica/solar) están generando una demanda sin precedentes de cobre, litio, níquel y cobalto. Los data centers de IA añaden una capa adicional de demanda intensiva de cobre de alta conductividad.

\subsubsection{Inversión en Seguridad de Suministro}
Gobiernos de EE.UU., UE y aliados están destinando recursos masivos a la diversificación de cadenas de suministro de minerales críticos, reduciendo la dependencia de China (que controla 91\% del refino de tierras raras). Esto crea un mercado nuevo para estudios de factibilidad y diseño de plantas de procesamiento en jurisdicciones occidentales.

\subsubsection{Digitalización de Operaciones Mineras}
El mercado de Smart Mining crece de USD 17.2B (2026) a USD 31.4B (2034) — CAGR 7.81\%. Incluye IoT, automatización, vehículos autónomos, gemelos digitales y mantenimiento predictivo con IA. Genera demanda de consultores con capacidades de integración tecnológica.

\subsubsection{Complejidad Regulatoria y ESG}
Las exigencias crecientes de cumplimiento ambiental, social y de gobernanza incrementan la demanda de consultoría especializada. El costo de compliance ESG se ha convertido en un componente significativo del CAPEX de nuevos proyectos.

\subsection{Análisis PESTLE}

\begin{table}[htbp]
\centering
\caption{Análisis PESTLE: Factores Clave del Entorno}
\begin{tabular}{@{}p{2cm}p{5.5cm}p{2cm}p{2.5cm}@{}}
\toprule
\textbf{Factor} & \textbf{Descripción} & \textbf{Impacto} & \textbf{Tendencia} \\
\midrule
Político & Sanciones a Rusia; proteccionismo en EE.UU.; reformas de permisos en Chile/Perú & Crítico & \trendup{} Creciente \\
\rowcolor{tablealt} Económico & Precios del cobre en máximos; inflación en costos operativos; TC multi-moneda & Alto & \trendflat{} Estable \\
Social & Escasez de talento ingenieril; comunidades indígenas en LATAM; licencia social & Alto & \trendup{} Creciente \\
\rowcolor{tablealt} Tecnológico & IA/IoT en minería; automatización; energías renovables en operaciones & Positivo & \trendup{} Creciente \\
Legal & Regulación ambiental más estricta; cambios en regalías (Chile, Perú); compliance & Alto & \trendup{} Creciente \\
\rowcolor{tablealt} Ambiental & Escasez hídrica en norte de Chile; descarbonización obligatoria; cierre de minas & Crítico & \trendup{} Creciente \\
\bottomrule
\end{tabular}
\label{tab:pestle}
\end{table}

\subsection{Tendencias Emergentes}

\begin{enumerate}
    \item \textbf{Minería con energía renovable:} Operaciones 100\% solar/eólica en el desierto de Atacama y Pilbara. Las consultoras que integren diseño energético con planificación minera tendrán ventaja.
    \item \textbf{Reciclaje y economía circular:} El cobre reciclado podría proveer hasta 25\% de la oferta necesaria para 2035. Nuevo nicho de consultoría en diseño de plantas de reciclaje.
    \item \textbf{Procesamiento doméstico (nearshoring):} EE.UU.\ construye capacidad de refino en instalaciones militares. LATAM busca agregar valor local vs.\ exportar concentrado.
    \item \textbf{Leyes de mineral decrecientes:} Las minas existentes producen menos cobre por tonelada procesada, lo que incrementa la intensidad energética, hídrica y de ingeniería — beneficiando a las consultoras.
\end{enumerate}
